\section{Заключение}
\label{sec:Chapter6} \index{Chapter6}

В ходе работы удалось достичь всех заявленных целей в достаточно полном объеме, а именно:

\begin{itemize}
\item изучена предметная область~--- проанализированы устройство OCI-образов, особенности их хранения и распространения через container registry, а также ключевые принципы работы современных сканеров уязвимостей. Сформировано понимание того, какая часть содержимого образа фактически нужна сканеру при индексировании;

\item произведён обзор существующих сканеров уязвимостей с открытым исходным кодом~--- Trivy, Clair, Grype. Изучены их архитектуры и модели работы с container registry, в результате чего выявлен общий недостаток: при сканировании образа из registry необходимые для индексирования слои скачиваются целиком, в то время как для анализа в большинстве случаев нужна лишь малая часть их содержимого;

\item предложен подход к решению выявленной проблемы, основанный на использовании формата eStargz и ленивой загрузке файлов из слоев через TOC. Подход обладает рядом важных свойств: совместимостью с существующими спецификациями OCI, отсутствием изменений на стороне registry, обратной совместимостью со стандартными container runtime;

\item разработан прототип сканера на базе Trivy, реализующий предложенный подход;

\item разработан фреймворк для воспроизводимого бенчмарка, моделирующий реальный профиль нагрузки на container registry;

\item проведена экспериментальная оценка разработанного прототипа, по результатам которой подтверждена применимость предложенного подхода. Достигнуты следующие количественные показатели:
\begin{itemize}
\item кратное (4--10X) ускорение сканирования по сравнению с обычным Trivy;
\item снижение объёма скачиваемых из registry данных до 5--15\% от размера образа в большинстве случаев;
\item совпадение результатов сканирования с обычным Trivy на всех образах тестового набора (с учётом оговорки про opaque whiteout, обсуждавшейся в \hyperref[sec:Chapter4]{разделе 5});
\end{itemize}

\end{itemize}

\subsection{Ограничения подхода и направления дальнейшего развития}

Несмотря на достигнутые результаты, предложенный подход обладает рядом ограничений, которые имеет смысл обсудить отдельно. Эти ограничения определяют возможные направления его дальнейшего развития:

\begin{itemize}
\item \textit{Зависимость от формата хранения.} Предложенный подход требует, чтобы образы хранились в registry в формате eStargz. На стороне пользователя это означает дополнительный шаг конвертации в CI/CD-конвейере, либо настройку сборщика на немедленное создание eStargz-варианта. Стоит, однако, отметить, что эта необходимость не является фундаментальной: как обсуждалось в \hyperref[sec:Chapter4]{разделе 5}, ряд managed container registry (Azure Container Registry в режиме Artifact Streaming~\cite{acr_artifact_streaming}, Amazon ECR с SOCI Index Builder~\cite{ecr_soci_builder}) уже самостоятельно генерируют соответствующий артефакт сразу после загрузки образа, не требуя от пользователя никаких дополнительных действий. По мере распространения подобных механизмов это ограничение перестанет быть существенным.

\item \textit{Метаданные, встроенные в исполняемый файл.} Как было показано в экспериментальной оценке, для образов формата \texttt{scratch} с Go-бинарным файлом ленивая загрузка на уровне отдельных файлов слоя не даёт выигрыша, поскольку метаинформация о модулях вшита непосредственно в исполняемый файл, и анализатору приходится запрашивать его целиком. Содержательно эту неэффективность можно было бы преодолеть, если бы Trivy умел читать только нужные фрагменты файла~--- например, секцию \texttt{.go.buildinfo} в ELF-файле, по которой и определяется состав модулей. Однако в текущем виде формат eStargz этого не позволяет: каждый файл упакован в отдельный gzip-поток, и считать произвольный фрагмент файла отдельным Range-запросом нельзя~--- gzip-поток нужно прочитать с начала. Поэтому обеспечение частичной загрузки фрагментов файлов в условиях пофайлового сжатия представляется отдельным интересным направлением для дальнейшего исследования.

\item \textit{Обработка opaque whiteout.} Текущая реализация прототипа не учитывает opaque whiteout, что в общем случае может приводить к погрешности в результатах сканирования. На практике в проведённом эксперименте такие случаи не встречались, однако в production-реализации эту доработку следует выполнить.

\item \textit{Дедупликация данных между слоями.} Использование eStargz приводит к небольшому увеличению размера образа в registry за счёт перехода на пофайловое сжатие. В качестве естественного направления развития интересно рассмотреть подход, применяемый в формате Nydus~\cite{nydus}: данные слоя разбиваются на \textit{chunk'и}, которые адресуются по содержимому и переиспользуются между разными слоями и образами. На наборах образов с большим количеством общих файлов такая дедупликация позволяет существенно сократить совокупный объём хранения и при достаточной её эффективности~--- полностью нивелировать overhead, добавляемый пофайловым сжатием.

\end{itemize}

\subsection{Итог}

Подытоживая, можно сказать, что данная работа предлагает практичное решение проблемы избыточной загрузки данных при сканировании OCI-образов из container registry. Решение основано на использовании существующего формата eStargz, разработанного в рамках развития ленивой загрузки контейнерных образов, и адаптации одного из наиболее популярных сканеров уязвимостей под этот формат. Реализованный прототип демонстрирует существенные количественные показатели~--- кратное ускорение и многократное сокращение трафика. Полученные результаты позволяют рекомендовать предложенный подход как направление развития существующих сканеров уязвимостей.

Все наработки, полученные в ходе работы, доступны в открытом доступе:
\begin{itemize}
\item исходный код форка Trivy с поддержкой eStargz~--- \url{https://github.com/loochek/trivy/tree/estargz};
\item фреймворк для воспроизводимого бенчмарка~--- \url{https://github.com/loochek/msc-work}.
\end{itemize}

\newpage
